\section{Background and Predictions}

\subsection{Client-centric consistency}

Let \(v\) denote an integer application version allocated by the single
logical writer. A write is \(W_c(x,v)\) and a read is \(R_c(x,v)\). The
predicates compare versions and dependencies observed by the client rather
than wall-clock timestamps.

\paragraph{Course definitions.}

\noindent
The following statements use the definitions given in Lecture~3. They are
separate from the project-specific checks described below.

\begin{description}
\item[RYW] The effect of a write operation by a process on data item \(x\) will
always be seen by a successive read operation on \(x\) by the same process
\cite[slides 37 to 38]{cai-lecture-3}.
\item[MR] If a process reads the value of a data item \(x\), any successive
read operation on \(x\) by that process will always return that same value or a
more recent value \cite[slides 41 to 42]{cai-lecture-3}.
\item[MW] A write operation by a process on a data item \(x\) is completed
before any successive write operation on \(x\) by the same process
\cite[slides 43 to 44]{cai-lecture-3}.
\item[WFR] A write operation by a process on a data item \(x\) following a
previous read operation on \(x\) by the same process is guaranteed to take
place on the same or a more recent value of \(x\) that was read
\cite[slides 39 to 40]{cai-lecture-3}.
\end{description}

\paragraph{Operationalisation in this project.}
RYW and MR are tested directly from successive versions returned to the
test client. The recorder also captures the evidence needed for the two
write properties.
For MW, the atomic pre-image of the successive write records the update
identifiers already present at the member executing that write. The checker
reports a violation when the preceding write identifier is absent from this
pre-image. For WFR, the same pre-image records the highest application version
present immediately before the dependent write as
\texttt{write\_base\_version}; the checker compares it with the version
returned by R1.

\subsection{MongoDB mechanisms}

We vary three mechanisms. A \texttt{local} read can see a member's current
state, including an update that has not been majority committed; a
\texttt{majority} read is limited to majority-committed data
\cite{mongodb-read-concern}. With \texttt{w:1}, a write can be acknowledged by
the primary before it reaches a majority. \texttt{majority} write concern
waits for majority acknowledgement \cite{mongodb-write-concern}. A causal
session carries ordering information between operations; MongoDB documents
its strongest causal guarantees for majority reads and majority writes
\cite{mongodb-causal-consistency}. We also choose where a read is attempted
through read preference and routing; the recorded history identifies the
member that handled it \cite{mongodb-read-preference}. These settings determine
which states may be returned, when writes are acknowledged, and which
dependencies later operations carry. We use these documented behaviours to
form the schedule-specific predictions below.

\subsection{Configuration matrix and predictions}

We combine the two read concerns (\texttt{local}, \texttt{majority}), two write
concerns (\texttt{w:1}, \texttt{majority}), and causal sessions switched off or
on. This produces eight configurations. For each property, we keep the
operation sequence, logical document, and replica-set topology fixed while
changing these controls. Table~\ref{tab:configurations} marks the
configuration-property pairs for which we made a positive prediction; the
other cells are descriptive comparisons, not claims of a guarantee.

\begin{table}[ht]
\centering
\small
\setlength{\tabcolsep}{3pt}
\caption{C1-C8 settings and properties expected under the target schedules.}
\label{tab:configurations}
\begin{tabular}{@{}l l l l p{0.30\linewidth}@{}}
\toprule
Config. & Read concern & Write concern & Causal session & Expected properties \\
\midrule
C1 & local & \texttt{w:1} & off & Descriptive only \\
C2 & local & \texttt{w:1} & on & Descriptive only \\
C3 & majority & \texttt{w:1} & on & MR, WFR \\
C4 & local & majority & on & MW \\
C5 & majority & majority & off & Descriptive only \\
C6 & majority & majority & on & RYW, MR, MW, WFR \\
C7 & majority & \texttt{w:1} & off & Descriptive only \\
C8 & local & majority & off & Descriptive only \\
\bottomrule
\end{tabular}
\end{table}

C1 leaves reads local, acknowledges one replica, and carries no causal
dependency; we therefore expect a stale or reordered result to be possible
when the schedule splits the replicas. C3 asks whether majority reads and a
causal session preserve MR and WFR in the tested sequence even though its
write uses \texttt{w:1}; this is a schedule-specific prediction, not a general
MongoDB guarantee for that combination. C4 tests whether an acknowledged
majority write preserves MW when a successor write follows it; its local read
does not protect WFR. C6 combines majority reads, majority writes, and a
causal session. We expect its dependencies not to produce a completed
counterexample, although a partition may prevent an operation from finishing.

We compare these predictions only with completed, decidable histories. A
timeout or missing precondition leaves the prediction untested in that trial.
In particular, a stale RYW read in C7 describes an unprotected setting; it
does not falsify a guarantee we assigned to C7.

\subsection{Fault expectations}

The first campaign changes settings while its schedule creates controlled
replica divergence. The second asks how the same client dependencies behave
when we change the topology.

\begin{description}
\item[F1, secondary loss] Removing one secondary should leave the primary and
a majority side intact. We expect the scheduled calls to complete without a
new consistency counterexample.
\item[F2, primary loss] Stopping the primary should trigger an election. We
expect calls made after the election barrier to preserve the tested
dependencies, but we make no prediction about a call during the election
because our schedule does not issue one.
\item[F3, replication partition] Isolating the former primary should separate
a \texttt{w:1} acknowledgement from a majority acknowledgement. We predict
that weak settings may complete with stale or reordered state, whereas C6 may
wait or leave a dependent operation unresolved.
\end{description}

